\uuid{BiXD}
\titre{Solution particulière/condition initiale}
\auteur{hanani ; mzari}
\organisation{exo7}
\licence{CC BY-NC-SA 4.0}
\datecreate{2020-01-20}
\niveau{L1}
\module{Analyse}
\chapitre{Equation différentielle}
\sousChapitre{Résolution d'équation différentielle du premier ordre}
\difficulte{3}
\theme{Equations du premier ordre, solution particulière, condition initiale}
\exo7qcmid{lille-2.equations-differentielles.013}

\contenu{
\begin{question}
\qclassification{225.01}
\qtags{motcle=solution particulière/condition initiale}
On considère l'équation différentielle $(E)$ : $\displaystyle x^2y'-y=1$ sur $\Rr$. Parmi les affirmations suivantes cocher celles qui sont vraies :
\begin{answers}
\good{La fonction $\displaystyle y=-1+\mathrm{e}^{-1/x}$ est une solution de $(E)$ sur $]0,+\infty[$.}
\bad{La fonction $\displaystyle y=1-\mathrm{e}^{-1/x}$ est une solution de $(E)$ sur $]-\infty,0[$.}
\bad{La fonction $\displaystyle y=2\mathrm{e}^{-1/x}$ est une solution de $(E)$ sur $\Rr$.}
\good{La solution de $(E)$ sur $]0,+\infty[$ vérifiant $y(1)=-1$ est constante.}
\end{answers}
\begin{explanations}
La fonction $y=-1+\mathrm{e}^{-1/x}$ est une solution de $(E)$ sur $]0,+\infty[$. Par contre $\displaystyle y=1-\mathrm{e}^{-1/x}$ n'est pas une solution de $(E)$ sur $]-\infty,0[$. Toute solution de $(E)$ sur $\Rr$ vérifie $y(0)=-1$, donc $\displaystyle y(x)=2\mathrm{e}^{-1/x}$ ne peut être une solution de $(E)$ sur $\Rr$. La solution générale de $(E)$ sur $]0,+\infty[$ est $y=-1+k\mathrm{e}^{-1/x}$, $k\in \Rr$. Donc, si $y(1)=-1$, alors $k=0$.
\end{explanations}
\end{question}
}
